Data centres' growing electricity footprint has motivated a body of
research on carbon-aware scheduling, in which workloads are shifted in
time or space to coincide with cleaner electricity. Much of the
foundational carbon-aware scheduling literature, however, targets
batch or long-running jobs, while recent FaaS-specific systems each
address only part of the spatial, temporal, hardware, and autoscaling
design space. The assumptions of the batch-oriented systems
--- minutes-to-hours runtime, delay tolerance, and clean
suspend/resume --- are profoundly violated by serverless or
Function-as-a-Service (FaaS) workloads. FaaS functions execute in
milliseconds to seconds, often under sub-second SLAs, suffer
expensive cold starts when interrupted, and arrive in bursty,
event-driven patterns.

We present \GreenFaaS, a carbon-aware scheduling framework for
serverless workloads that unifies three decisions previously treated
separately: spatial routing across regions, temporal deferral within
SLA bounds, and warm-pool-aware container reuse. Four prior FaaS-targeted
carbon-aware schedulers cover different facets of the design space:
\emph{GreenCourier} \citep{chadha2023greencourier} routes single
functions to low-carbon regions; \emph{Caribou} \citep{gsteiger2024caribou}
shifts whole serverless workflows geographically; \emph{EcoLife}
\citep{jiang2024ecolife} exploits multi-generation hardware; and
\emph{CASA} \citep{qi2024casa} co-optimizes autoscaling and SLOs within
a single cluster. \GreenFaaS{} is the first to give a closed-form
joint analysis of the spatial and temporal axes under cold-start
economics, and to demonstrate empirically-validated robustness across topology and
workload regimes.

The paper's principal analytical contribution is a closed-form
\emph{cold-start carbon trade-off}: we prove a dichotomy in which the
warm-pool placement decision is governed by a dimensionless idle-cost
ratio $\beta = P_i T_w / (P_a \tau_e)$ relative to a critical
threshold $(r-1)(1+\alpha)$, with a closed-form break-even condition and a unique critical
arrival rate $\lambda^*$ --- the root of a monotone scalar equation
--- separating the regimes. The same idle-cost ratio governs
the analogous temporal-shift trade-off with the deferral interval
$\Delta t$ in place of $T_w$, giving the scheduler principled
decision rules along both axes.

We evaluate \GreenFaaS{} in a discrete-event simulator, driven by a
schema-faithful synthetic workload calibrated to the Azure Functions
2019 and 2021 characterizations and by real grid carbon-intensity
data from Let's-Wait-Awhile \citep{wiesner2021lets,lwa} and from
ElectricityMaps over the same January 2021 window as the Azure trace,
against four carbon-aware baselines representative of recent work.
Three findings stand out.
\emph{First}, \GreenFaaS{} reduces operational carbon by 68--83\% relative to FIFO across
a range of topologies on fully-real time-aligned data (real Azure
2021 trace + real ElectricityMaps Jan 2021 carbon), and 64--72\% on
synthetic-workload setups. On fully-real time-aligned data the
\GreenFaaS{}--Spatial gap is 0.04 percentage points; across all
topologies on real-carbon data the gap stays within 0.5pp.
\GreenFaaS's distinct value over Spatial is topology-robustness
rather than dominance in any single regime.
\emph{Second}, \GreenFaaS{} preserves a \emph{do-no-harm} property
that extends beyond zero-opportunity workloads to a robust
neighbourhood of low-opportunity ones: across shiftable fractions
$\{0, 1\%, 2\%\}$ and in single-region scenarios on both synthetic
and fully-real data, \GreenFaaS{} matches FIFO byte-for-byte; a
heuristic ablation lacking the cold-start trade-off correction loses
5--102\% to FIFO in the same scenarios.
\emph{Third}, the canonical Wait-Awhile temporal scheduler and a
FaaS-adapted port of the provably-optimal one-way trading algorithm
of \citet{lechowicz2024online} both \emph{fail} on FaaS workloads,
losing 3--306\% to FIFO depending on the regime --- catastrophically
(up to $-306\%$) at batch-default deferral horizons on synthetic
workloads, and mildly (3--5\%) on fully-real data where they
degenerate toward a FIFO no-op. In neither regime do they reduce carbon: a clean
failure mode for batch carbon-aware techniques on FaaS, including
those with provable competitive guarantees, when warm-pool
idle-energy coupling is not accounted for. Head-to-head against the closest competing
serverless system, \emph{Caribou} \citep{gsteiger2024caribou}, well-tuned
Caribou (hourly re-deployment with perfect carbon forecasting) and
\GreenFaaS{} achieve statistically indistinguishable carbon savings (within 0.1pp) on real Azure
2021 workload; the difference is operational rather than empirical,
with \GreenFaaS{} avoiding the container migration and pub/sub
orchestration that Caribou requires.
The simulator, scheduler, trace loaders, Caribou port, and
real Azure Functions and ElectricityMaps integrations are released as
open-source artifacts under the Apache-2.0 license upon publication.
